\documentclass[12pt,letterpaper]{article}
\usepackage{fontspec}
\usepackage[spanish,es-noshorthands,es-tabla,es-lcroman]{babel}
\usepackage{../plantilla/latex/unab_reporte}

% Los datos compartidos se generan desde datos/datos.md al ejecutar compilar.py.
\input{../.build/datos.tex}
\input{../.build/citas.tex}
\hypersetup{pdftitle={\TituloTrabajo},pdfauthor={\NombreFirma}}

\begin{document}
\portadareporte
\pagenumbering{roman}
\section*{Resumen}
\addcontentsline{toc}{section}{Resumen}
Escribe aquí una síntesis del reporte: el problema abordado, el objetivo,
la metodología utilizada, los principales resultados y la conclusión central.
Redacta este apartado cuando hayas terminado el desarrollo del documento.

\textbf{Palabras clave:} tema; concepto; área de estudio.
\clearpage
\tableofcontents
\clearpage
\pagenumbering{arabic}

\section{Introducción}
Presenta el tema, su contexto y su importancia. Describe el problema o la
pregunta que motiva el trabajo, el alcance del reporte y cómo está organizado.
Reemplaza todos los textos de orientación con el contenido de tu investigación.

\section{Objetivos}
\subsection{Objetivo general}
Expresa el propósito principal del trabajo con un verbo en infinitivo.
\subsection{Objetivos específicos}
\begin{itemize}
  \item Identificar los elementos principales del problema de estudio.
  \item Analizar la información recopilada mediante criterios definidos.
  \item Proponer recomendaciones a partir de los resultados obtenidos.
\end{itemize}

\section{Marco teórico}
\subsection{Conceptos fundamentales}
Desarrolla las definiciones y los antecedentes necesarios para comprender el
tema. Respalda las afirmaciones con fuentes y citas, por ejemplo~\cita{fuente}.
La referencia incluida al final es un modelo que debes reemplazar.
\subsection{Antecedentes}
Describe los trabajos previos y explica su relación con el problema estudiado.

\section{Metodología}
Explica el enfoque del trabajo, las fuentes de información, los instrumentos
utilizados y el procedimiento seguido. Indica los criterios de selección y
las limitaciones para que el proceso pueda comprenderse y reproducirse.

\section{Resultados y análisis}
\subsection{Presentación de resultados}
Expón los hallazgos de forma ordenada. La tabla~\ref{tab:resultados} muestra
una estructura de ejemplo; sustituye las filas con tus datos.
\begin{table}[htbp]
  \centering
  \caption{Síntesis de resultados (ejemplo de estructura)}
  \label{tab:resultados}
  \begin{tabularx}{\linewidth}{@{}lXX@{}}
    \toprule
    \textbf{Criterio} & \textbf{Resultado} & \textbf{Interpretación} \\
    \midrule
    Criterio 1 & Describe el hallazgo. & Explica su significado. \\
    Criterio 2 & Describe el hallazgo. & Relación con el objetivo. \\
    Criterio 3 & Describe el hallazgo. & Implicaciones principales. \\
    \bottomrule
  \end{tabularx}
  \par\smallskip {\small Fuente: elaboración propia.}
\end{table}

\subsection{Discusión}
Interpreta los resultados y contrástalos con el marco teórico. Distingue entre
los datos observados y tus interpretaciones. Explica las limitaciones del análisis.

% Para destacar una nota u observación importante, usa:
% \begin{cajanota}[Nota metodológica]
%   Escribe aquí el texto destacado, hallazgo o consideración relevante.
% \end{cajanota}

% Para insertar un bloque de código fuente con estilo UNAB, descomenta:
% \begin{lstlisting}[language=Python, caption={Ejemplo de procesamiento de datos}]
% def procesar_datos(registros):
%     return [item.strip() for item in registros if item]
% \end{lstlisting}

% Para insertar una figura, guarda el archivo en assets/contenido/imagenes/ y descomenta:
% \begin{figure}[htbp]
%   \centering
%   \includegraphics[width=0.8\linewidth]{../assets/contenido/imagenes/resultado.png}
%   \caption{Descripción de la figura. Fuente: elaboración propia.}
%   \label{fig:resultado}
% \end{figure}

\section{Conclusiones}
Responde al objetivo general y sintetiza los hallazgos principales. Explica
qué se puede concluir a partir de la evidencia presentada, sin introducir
resultados nuevos.

\section{Recomendaciones}
Propón acciones concretas derivadas de las conclusiones. Señala responsables,
condiciones o líneas de investigación futura cuando corresponda.

\clearpage
\phantomsection
\addcontentsline{toc}{section}{Referencias}
\section*{Referencias}
\begin{doublespace}
\input{../.build/referencias-reporte.tex}
\end{doublespace}

\clearpage
\appendix
\section{Material complementario}
Incluye instrumentos, datos adicionales, evidencias u otros materiales que
complementen el reporte. Elimina esta sección si no necesitas anexos.

% ==============================================================================
% FIRMA INSTITUCIONAL: Se imprime obligatoriamente al final de todo el reporte,
% después de los anexos (o tras las referencias si no hay anexos).
% Consérvala siempre en el documento final.
% ==============================================================================
\firmareporte
\end{document}
